\documentclass[12pt]{article}
\usepackage[a4paper, margin=1in]{geometry}
\usepackage{mathptmx} % Times New Roman
\usepackage{amsmath, amssymb}
\usepackage{graphicx}
\usepackage{booktabs}
\usepackage{array}
\usepackage{enumitem}
\usepackage{hyperref}

\hypersetup{
    colorlinks=false,
    pdfborder={0 0 0}
}

\setlist{itemsep=2pt, topsep=3pt}

\begin{document}

\begin{center}
    {\Large \textbf{Experiment 3}}
\end{center}

\vspace{0.5em}

\noindent \textbf{Aim:} Exploratory Data Analysis \& Statistical Analysis

\vspace{0.8em}

\noindent \textbf{Objective:}
\begin{enumerate}[leftmargin=1.8em]
    \item To visualize the distribution of classes and features in the dataset.
    \item To understand the spread and central tendency of data using plots.
    \item To identify correlations between features using heatmaps.
    \item To perform statistical hypothesis testing (e.g., t-tests, ANOVA, Chi-Square) to determine if observed differences are significant (as per the requirement).
\end{enumerate}

\vspace{0.8em}

\noindent \textbf{Dataset Overview \& Summary Statistics}

\noindent The analysis is performed on the Customer Support on Twitter (TWCS) Cleaned Dataset (\texttt{twcs\_cleaned.csv}) containing 100,000 processed customer support tweets downsampled and stratified for robust statistical evaluation.

\vspace{0.5em}

\begin{table}[h!]
\centering
\small
\begin{tabular}{lrrrrrrr}
\toprule
\textbf{Feature} & \textbf{Mean} & \textbf{Std Dev} & \textbf{Min} & \textbf{25\% (Q1)} & \textbf{50\% (Median)} & \textbf{75\% (Q3)} & \textbf{Max} \\
\midrule
Word Count ($w$)        & 19.23  & 9.68  & 5.00   & 12.00  & 18.00  & 24.00  & 68.00  \\
Character Count ($c$)   & 102.23 & 51.34 & 11.00  & 64.00  & 100.00 & 127.00 & 280.00 \\
VADER Compound          & +0.028 & 0.456 & -0.991 & -0.309 & 0.000  & +0.402 & +0.984 \\
Positive Valence ($p$)  & 0.103  & 0.132 & 0.000  & 0.000  & 0.060  & 0.172  & 0.919  \\
Negative Valence ($n$)  & 0.085  & 0.113 & 0.000  & 0.000  & 0.000  & 0.147  & 0.943  \\
Neutral Valence ($u$)   & 0.812  & 0.154 & 0.057  & 0.713  & 0.824  & 0.941  & 1.000  \\
Exclamation Count ($!$) & 0.461  & 1.042 & 0.000  & 0.000  & 0.000  & 1.000  & 24.00  \\
Question Count ($?$)    & 0.528  & 0.793 & 0.000  & 0.000  & 0.000  & 1.000  & 11.00  \\
\bottomrule
\end{tabular}
\end{table}

\vspace{0.8em}

\noindent \textbf{Detailed Steps}

\begin{itemize}[leftmargin=1.5em]
    \item \textbf{Plot class balance:} Visualized the frequency distribution across the 5 target emotion classes using count bar plots and proportional pie charts.
    \begin{itemize}
        \item Neutral / Inquiry: 20,111 samples (40.22\%)
        \item Disappointment / Sadness: 13,091 samples (26.18\%)
        \item Joy / Gratitude: 10,701 samples (21.40\%)
        \item Anger / Frustration: 5,385 samples (10.77\%)
        \item Fear / Anxiety: 712 samples (1.42\%)
    \end{itemize}

    \item \textbf{Visualize frequency distributions of features:} Evaluated the shape, central tendency, and dispersion of continuous features (\texttt{word\_count}, \texttt{char\_count}, \texttt{vader\_compound}, \texttt{exclamation\_count}) using Histograms with Kernel Density Estimation (KDE), parametric mean indicators, and non-parametric median indicators.

    \item \textbf{Assess Feature Correlations:} Computed Pearson ($r$) linear correlation and Spearman ($\rho$) monotonic rank correlation matrices. Visualized lower-triangle annotated correlation heatmaps to assess feature collinearties.

    \item \textbf{Feature Distribution Analysis \& Outlier Detection:} Fitted theoretical continuous and discrete probability distributions (Gaussian, Log-Normal, Exponential, and Poisson) to tweet word counts using maximum likelihood estimation. Evaluated goodness-of-fit via Kolmogorov-Smirnov (K-S) tests and Normal Q-Q plots. Applied Tukey's Fences (IQR method) and Z-score ($|Z| > 3$) outlier detection.

    \item \textbf{Hypothesis Testing:} Formulated and executed statistical significance tests:
    \begin{itemize}
        \item Independent Two-Sample Welch's t-test (Negative complaint vs.\ Joyful praise word counts)
        \item One-Way Analysis of Variance (ANOVA) and Tukey's Honestly Significant Difference (HSD) post-hoc test (VADER polarity across emotion classes)
        \item Chi-Square ($\chi^2$) Test of Independence (Emotion distribution vs.\ Time of Day)
        \item Mann-Whitney U Test (Non-parametric rank sum comparison of Anger vs.\ Neutral word counts)
    \end{itemize}

    \item \textbf{Document EDA insights:} Documented statistical inferences, distributional topologies, skewness, and domain findings.
\end{itemize}

\vspace{0.8em}

\noindent \textbf{Open-Source Tools}

\noindent Matplotlib, Seaborn, Plotly, SciPy, Statsmodels, Pandas, NumPy

\vspace{0.8em}

\noindent \textbf{Deliverables}

\noindent A well-commented EDA Notebook and Pipeline including:
\begin{enumerate}[leftmargin=1.8em]
    \item Data loading (\texttt{twcs\_cleaned.csv}) and multidimensional feature extraction
    \item Visualizations (class balance charts, frequency distributions, boxplots, violin plots, and correlation heatmaps)
    \item Statistical test implementations (Welch's t-test, ANOVA, $\chi^2$ independence, Mann-Whitney U, and K-S distribution fitting)
    \item Summary of insights and empirical conclusions
\end{enumerate}

\newpage

\noindent \textbf{Hypothesis test results, clearly documenting:}

\vspace{0.5em}

\noindent \textbf{1. Two-Sample Welch's t-Test (Negative Complaint vs. Joy Word Count)}
\begin{itemize}[leftmargin=1.5em]
    \item \textbf{Hypotheses:}
    \begin{align*}
        H_0 &: \mu_{\text{negative}} = \mu_{\text{joy}} \quad \text{(Mean word count of complaints equals that of joyful praise)} \\
        H_1 &: \mu_{\text{negative}} \neq \mu_{\text{joy}} \quad \text{(Mean word count differs significantly)}
    \end{align*}
    \item \textbf{Test Statistic \& Parameters:}
    \begin{itemize}
        \item Negative Cohort ($N_1 = 18,476$): $\bar{x}_1 = 21.29$ words, $s_1 = 9.91$
        \item Joy Cohort ($N_2 = 10,701$): $\bar{x}_2 = 16.78$ words, $s_2 = 9.13$
        \item Degrees of Freedom: $df = 23,892.4$
        \item Test Statistic: $t = 39.3351$
        \item Effect Size (Cohen's $d$): $d = 0.468$ (Moderate effect size)
    \end{itemize}
    \item \textbf{P-value and interpretation:} $p = 0.0000$ ($p < 10^{-10}$, $\alpha = 0.05$). Since $p < \alpha$, the probability of observing this mean difference under $H_0$ is negligible.
    \item \textbf{Conclusion based on test result:} \textbf{Reject $H_0$}. Customers filing complaints express significantly greater verbosity ($+4.51$ words on average, a $26.9\%$ increase) than customers offering praise.
\end{itemize}

\vspace{0.6em}

\noindent \textbf{2. One-Way Analysis of Variance (ANOVA) \& Tukey HSD (Emotion Polarity Separation)}
\begin{itemize}[leftmargin=1.5em]
    \item \textbf{Hypotheses:}
    \begin{align*}
        H_0 &: \mu_{\text{anger}} = \mu_{\text{sadness}} = \mu_{\text{fear}} = \mu_{\text{neutral}} = \mu_{\text{joy}} \\
        H_1 &: \text{At least one emotion category has a significantly different mean polarity.}
    \end{align*}
    \item \textbf{Group Means (VADER Compound):} Anger: $-0.5543$, Sadness: $-0.3550$, Fear: $-0.1486$, Neutral: $+0.1356$, Joy: $+0.6000$.
    \item \textbf{Test Statistic \& P-value:} $F = 32,429.39$, $p = 0.0000$ ($p < 10^{-15}$, $\alpha = 0.05$).
    \item \textbf{Interpretation \& Post-Hoc Tukey HSD:} All pairwise between-group differences are statistically significant ($p < 0.001$).
    \item \textbf{Conclusion based on test result:} \textbf{Reject $H_0$}. The five emotion categories form distinct, non-overlapping sentiment polarity strata.
\end{itemize}

\vspace{0.6em}

\noindent \textbf{3. Chi-Square ($\chi^2$) Test of Independence (Emotion vs. Time-of-Day)}
\begin{itemize}[leftmargin=1.5em]
    \item \textbf{Hypotheses:}
    \begin{align*}
        H_0 &: \text{Emotion category is independent of Time of Day (Morning, Afternoon, Evening, Night)} \\
        H_1 &: \text{Emotion category distribution depends significantly on Time of Day}
    \end{align*}
    \item \textbf{Test Statistic \& Parameters:} $\chi^2 = 56.5628$, $df = 12$, Cram\'{e}r's $V = 0.0194$.
    \item \textbf{P-value and interpretation:} $p = 9.4784 \times 10^{-8}$ ($p < 10^{-5}$, $\alpha = 0.05$).
    \item \textbf{Conclusion based on test result:} \textbf{Reject $H_0$}. There is a statistically significant temporal relationship; angry inquiries surge proportionally during Evening and Night hours.
\end{itemize}

\vspace{0.6em}

\noindent \textbf{4. Theoretical Distribution Fitting \& Outlier Analysis Summary}
\begin{itemize}[leftmargin=1.5em]
    \item \textbf{Goodness-of-Fit (Kolmogorov-Smirnov Test):}
    \begin{itemize}
        \item Log-Normal Fit: $\text{Shape} = 0.39, \text{Scale} = 22.63 \implies \text{KS-Stat} = 0.0665$ (\textbf{Best Continuous Fit})
        \item Gaussian (Normal) Fit: $\mu = 19.23, \sigma = 9.68 \implies \text{KS-Stat} = 0.0809$
        \item Exponential Fit: $\text{Scale} = 14.23 \implies \text{KS-Stat} = 0.1640$
        \item Poisson Fit (Discrete): $\lambda = 19.23$
    \end{itemize}
    \item \textbf{Outlier Detection (Tukey's IQR Fences):} Upper fence at $42.0$ words isolates $1,799$ outliers ($3.60\%$). These extreme length tweets are predominantly high-urgency negative customer complaints.
\end{itemize}

\vspace{0.8em}

\noindent \textbf{Conclusion}

\noindent Experiment 3 successfully accomplished the Exploratory Data Analysis and Statistical Hypothesis Testing objectives. The empirical data exhibits a right-skewed Log-Normal distribution for tweet lengths ($KS = 0.0665$). Formal hypothesis tests confirmed with high statistical significance ($p < 10^{-10}$) that dissatisfied customer complaints demand substantially higher word counts ($+4.51$ words, $t = 39.34$) and exhibit elevated late-night volumes ($\chi^2 = 56.56, p < 10^{-7}$). All experimental deliverables, plots, and code implementations have been validated and documented.

\end{document}
